% Volume II: Causal Explosion
% Part IV. Time, Delay, and the Maturation of Evidence
% Chapter 17: Evidentiary Time
% Spine: proof-spines/ch017-spine.md

\chapter{Evidentiary Time}
\label{ch:ch017}

This chapter defines \textbf{evidentiary time} by separating event time \(t_e\), recording time \(t_r\), and recognition time \(t_s\). A detail may be captured immediately and yet matter only much later:
\[
t_e\leq t_r\ll t_s.
\]
\Definition\ This temporal ordering fixes the chapter's vocabulary for evidentiary latency.
That gap is the \textbf{latency of evidence}. It is not an accidental inconvenience but a structural variable in inquiry, because what can responsibly be believed depends on where one stands relative to those three times.

Not every source of delay bears on that gap in the same way, so the Part fixes four kinds before treating any of them as a virtue. \textbf{Procedural delay} lets evidence mature between \(t_r\) and \(t_s\), and is the only kind inherently connected to epistemic restraint. \textbf{Capacity delay} arises because an institution lacks the investigators, instruments, or expertise to process what it already holds. \textbf{Strategic delay} is imposed by an actor who benefits when attention, witnesses, or resources dissipate before \(t_s\). \textbf{Pathological delay} results from institutional indifference or dysfunction with no compensating evidentiary work occurring in the interval. Whenever a later chapter credits delay with protecting judgment, the claim being made is about procedural delay specifically; the other three can destroy the very freedom this Part is trying to explain, and each should be named rather than left to hide inside a generic appeal to patience.

Evidence therefore has a temporal ecology. Traces appear, decay, corroborate, and change significance on different schedules, so a claim that is weakly warranted early can become strongly warranted later without the underlying event changing at all. The chapter treats delay as both a cost and a condition of truth, and asks how institutions should act when judgment must often be made before recognition time has caught up with recording time.
